\documentclass[a4paper]{article}
\usepackage[utf8]{inputenc}
\usepackage[russian,english]{babel}
\usepackage[T2A]{fontenc}
\usepackage[left=10mm, top=20mm, right=18mm, bottom=15mm, footskip=10mm]{geometry}
\usepackage{indentfirst}
\usepackage{amsmath,amssymb}
\usepackage{mathastext} 
\usepackage[italicdiff]{physics}
\usepackage{graphicx}
\usepackage{caption}
\usepackage{float}
\usepackage{caption}
\renewcommand{\thefootnote}{\fnsymbol{footnote}}
\usepackage{tablefootnote}
\usepackage{footmisc}
\usepackage{textcomp}
\usepackage{multicol}
\usepackage{multirow}
\usepackage[parfill]{parskip}
\usepackage{makecell}
\usepackage{array}  
\usepackage{siunitx}
\usepackage{booktabs}
\sisetup{
    exponent-product = \cdot,
    per-mode = symbol,
    separate-uncertainty = true,
    uncertainty-separator = \,
}
\renewcommand\cellalign{cc} 
\renewcommand\cellgape{\Gape[2pt]} 
\usepackage[utf8]{inputenc}\newcommand{\approxtext}[1]{\ensuremath{\stackrel{\text{#1}}{\approx}}}
\usepackage[colorlinks=true, urlcolor=blue, linkcolor=blue, citecolor=blue]{hyperref}
\graphicspath{{images/}}
\DeclareGraphicsExtensions{.pdf,.png,.jpg}
\usepackage{wrapfig}
\captionsetup{labelformat=empty}
\usepackage{caption}
\captionsetup[figure]{name=Рисунок}
\captionsetup[table]{name=Таблица}
%\usepackage{newtxtext,newtxmath} 
\title{\textbf{Отчет о выполнении лабораторной работы 4.3.3} \\[0.5em] \large\itshape Исследование разрешающей способности микроскопа методом Аббе}
\date{}
\author{Котляров Михаил, Б01-402}

\begin{document}

\maketitle
	
\section{Введение}

\subsubsection*{Цель работы:} определение дифракционного предела разрешения объектива микроскопа методом Аббе.


\subsubsection*{Оборудование и инструментальные погрешности}

\textbf{лазер (серийный номер HX214101032002}: длина волны $\lambda = 532$ нм\\
\textbf{кассета с набором сеток разного периода},\\
\textbf{линзы}: Л$_1$ ($F_1 = 110$ мм), Л$_2$ ($F_2 = 25$ мм),\\
\textbf{щель с микрометрическим винтом}: цена деления 0,01 мм,\\
\textbf{экран},\\
\textbf{линейка}: цена деления 0,5 мм.\\


\section{Экспериментальная установка}

\begin{figure}[H]
    \centering
    \includegraphics[width = 0.7\textwidth]{Pictures/setup.png}
    \caption{Рис. 1. Схема экспериментальной установки — модель проекционного
        микроскопа}
    \label{fig:setup}
\end{figure}

Схема установки приведена на Рис. \ref{fig:setup}. Предметом $P_{1}$ служат
сетки в кассете $C$. Линза Л$_1$ -- длиннофокусная, а Л$_2$ -- короткофокусная. В F
Устанавливаются диафрагмы D, с помощью сеток с разными периодами $d$ и щелевой
диафрагмы можно проверить третье соотношение. Период сеток может быть измерен
либо по расстоянию между дифракционными максимумами на экране, либо по
увеличенному с помощью микроскопа изображению сетки на экране. Пространственную
фильтрацию (получение наклонного изображение решётки) можно получить с помощью
подбора угла наклона и ширины вспомогательной щели.

\section{Выполнение}

\subsection{Определение периода решёток по пространственному спектру}

Проведём измерение параметров дифракционной картины для каждой из шести сеток. Закрепим кассету с сетками вблизи выходного окна лазера, включим излучение и на удалённом экране ($L = 100 \pm 1$~см) наблюдаем дифракционные картины. Для каждой сетки измерим расстояние $l$ между удалёнными максимумами одного направления (горизонтального или вертикального) и число промежутков $m$ между ними. Рассчитаем период каждой решётки по формуле для главных максимумов дифракции на двумерной решётке при малых углах:
\begin{equation}
    d = \frac{(m-1) \lambda L}{l},
    \label{eq:d_spectrum}
\end{equation}
где $m $ --- количество максимумов в отрезке длиной $l$ (по обе стороны от центрального максимума), $\lambda = 532$~нм --- длина волны лазера, $L = 100\pm1$~см --- расстояние от сетки до экрана. Полученные значения $d$ также занесём в таблицу~\ref{tab:periods_with_errors}. Эти значения будем считать эталонными при последующем сравнении методов.

% Таблица 1: Измерение периода решётки №1 по разным направлениям максимумов
\begin{table}[h!]
\centering
\begin{tabular}{|l|c|c|c|}
\hline
Направление измерения & $l_m$, мм & $m$ & $l_1 = l_m/m$, мм \\
\hline
по горизонтали & 215,5 & 5 & 53,875 \\
\hline
по вертикали & 215,5 & 5 & 53,875 \\
\hline
по диагонали & 228,5 & 4 & 53,858 \\
\hline
\end{tabular}
\caption{Таблица 1. Измерение расстояний между дифракционными максимумами для решётки №1}
\label{tab:grid1_directions}
\end{table}

По измерению первой решетки рассчитаем полную погрешность измерения с помощью линейки:\\ $\sigma_\text{случ} = 0,006$ мм, $\sigma_\text{сист} = 0,5$ мм, $\sigma = \sqrt{\left ( \sigma_\text{случ}\right )^2 + \left ( \sigma_\text{сист}\right )^2} = 0,5$ мм.
% Таблица 2: Расчёт периодов решёток с погрешностями
\begin{table}[h!]
\centering
\begin{tabular}{|c|c|c|c|c|}
\hline
№ сетки & $l$, мм & $m$ & $d$, мкм & $\varepsilon_d$, \% \\
\hline
\multicolumn{5}{|c|}{Верхние максимумы} \\
\hline
1 & 215,5 & 5 & $9{,}875 \pm 0{,}101$ & 1,03 \\
\hline
2 & 259,5 & 13 & $2{,}460 \pm 0{,}025$ & 1,02 \\
\hline
3 & 255,5 & 25 & $4{,}997 \pm 0{,}051$ & 1,02 \\
\hline
\multicolumn{5}{|c|}{Нижние максимумы} \\
\hline
1 & 215,5 & 5 & $9{,}875 \pm 0{,}101$ & 1,03 \\
\hline
2 & 256,0 & 13 & $24{,}94 \pm 0{,}025$ & 1,02 \\
\hline
3 & 232,5 & 23 & $50{,}34 \pm 0{,}051$ & 1,02 \\
\hline
\end{tabular}
\caption{Таблица 2. Рассчитанные периоды решёток с оценкой погрешностей}
\label{tab:periods_with_errors}
\end{table}



\subsection{Определение периода по увеличенному изображению (модель микроскопа)}

Соберём модель проекционного микроскопа. Отцентрируем систему так, чтобы пучок после прохождения линз попадал в ту же точку экрана. Расположим длиннофокусную линзу Л$_1$ ($F_1 = 110$~мм) на расстоянии $a_1$, чуть большем фокусного, от сетки; с помощью линзы Л$_2$ получим резкое увеличенное изображение сетки на экране. Для исключения влияния саморепродукции сфокусируемся на изображении тонкой проволочки, нанесённой на одну из сеток.

Измерим геометрические параметры системы: $a_1$, $b_1$ (расстояния от сетки до Л$_1$ и от Л$_1$ до промежуточного изображения), $a_2$, $b_2$ (от Л$_2$ до промежуточного изображения и до экрана). Измерим также суммарное расстояние $b_1 + a_2$ (``длина тубуса''). Из формулы тонкой линзы $a_{2}=\frac{b_{2} F_{2}}{b_{2}-F_{2}}$. Увеличение микроскопа $\Gamma=\frac{b_{1} b_{2}}{a_{1} a_{2}}$. Для каждой сетки измерим период увеличенного изображения на экране $d_{\text{экрана}}$ (расстояние между соседними линиями). Рассчитаем действительный период сетки:
\begin{equation}
    d_{\text{микро}} = \frac{d_{\text{экрана}}}{\Gamma}.
    \label{eq:d_micro}
\end{equation}

% Таблица 1: Геометрические параметры модели микроскопа
\begin{table}[h!]
\centering
\begin{tabular}{|c|c|c|}
\hline
Параметр & Значение \\
\hline
$b_2$, мм & 479 \\
\hline
$a_2 + b_1$, мм & 77,6  \\
\hline
$a_1$, мм & 135  \\
\hline
$a_2$, мм & 26,38 \\
\hline
$\Gamma$ & $6{,}89 \pm 0{,}07$ \\
\hline
\end{tabular}
\caption{Таблица 3. Геометрические параметры модели микроскопа и расчёт увеличения}
\label{tab:geometry_magnification}
\end{table}

% Таблица 2: Периоды решёток по увеличенному изображению
\begin{table}[h!]
\centering
\begin{tabular}{|c|c|c|c|c|}
\hline
№ сетки & $l$, мм & $m$ & $d$, мкм & $\varepsilon_d$, \% \\
\hline
\multicolumn{5}{|c|}{Верхние максимумы} \\
\hline 
1 & 268,0 & 39 & $10{,}95 \pm 0{,}016$ & 1,46 \\
\hline
2 & 95,0 & 27 & $21{,}13 \pm 0{,}033$ & 1,54 \\
\hline
3 & 30,0 & 20 & $48{,}90 \pm 0{,}108$ & 2,21 \\
\hline
\end{tabular}
\caption{Таблица 4. Периоды решёток, рассчитанные по увеличенному изображению}
\label{tab:periods_microscope}
\end{table}

\subsection{Определение периода по разрешающей способности (метод Аббе)}

Поместим щелевую диафрагму с микрометрическим винтом в заднюю фокальную плоскость~$F$ линзы~$Л_1$ (объектива модели микроскопа). Для каждой сетки будем уменьшать ширину щели~$D$, наблюдая за вторичным изображением на экране. Зафиксируем минимальное значение~$D$, при котором ещё различимо периодическое изображение сетки, то есть щель пропускает дифракционные максимумы порядков $m = 0, \pm 1$ (при дальнейшем сужении щели изображение вырождается в одномерную структуру).

Для третей сетки измерили минимальную ширину щели $D = 2{,}17$~мм. Рассчитаем период решётки по формуле Аббе для когерентного освещения:
\begin{equation}
    d_{\text{Аббе}} = \frac{2 F_1 \lambda}{D} = 53{,}94 \pm 0{,}25~\text{мкм}, \quad \varepsilon_d = 0{,}46\,\%.,
    \label{eq:d_abbe_final}
\end{equation}

Для остальных сеток теоретически ожидаемые значения минимальной ширины щели~$D$, рассчитанные из формулы~\eqref{eq:d_abbe_final} с использованием эталонных периодов $d_{\text{спектр}}$, приведены в таблице~\ref{tab:abbe_expected}.

\begin{table}[h!]
\centering
\begin{tabular}{|c|c|c|c|}
\hline
№ сетки & $D_{\text{теор}}$, мм & $d_{\text{спектр}}$, м & $\varepsilon_d$, \% \\
\hline
1 & 11,85 & $1{,}217 \cdot 10^{-4}$ & 1,03 \\
\hline
2 & 4,76 & $4{,}845 \cdot 10^{-5}$ & 1,02 \\
\hline
\end{tabular}
\caption{Таблица 5. Ожидаемые значения ширины щели для проверки формулы Аббе}
\label{tab:abbe_expected}
\end{table}

Однако микрометрический винт используемой щелевой диафрагмы имеет максимальный ход $D_{\max} \approx 4$~мм, поэтому для сеток с большим периодом (№\,1 и частично №\,2) не удалось достичь требуемой ширины щели. В связи с этим измерения разрешающей способности были выполнены только для первой сетки, период которой достаточно мал, чтобы соответствующее значение~$D$ попало в диапазон регулировки диафрагмы.

\subsection{Пространственная фильтрация и мультиплицирование (качественные наблюдения)}

Проведём опыт по пространственной фильтрации. Подберём сетку со средним периодом и установим щель в фокальной плоскости $F$ так, чтобы она пропускала максимум нулевого порядка, но блокировала максимумы первого порядка в поперечном направлении. Поворачивая щель, получим изображения при трёх ориентациях:
\begin{itemize}
    \item \textit{Вертикальная щель}: пропускает максимумы $(0, m_y)$ --- на экране наблюдается решётка с горизонтальными штрихами;
    \item \textit{Горизонтальная щель}: пропускает максимумы $(m_x, 0)$ --- изображение с вертикальными штрихами;
    \item \textit{Щель под углом $45^\circ$}: пропускает максимумы с $m_x = m_y$ --- на экране возникает наклонная периодическая структура, отсутствующая в исходном объекте.
\end{itemize}
Наблюдаемые явления объясняются тем, что щель в фурье-плоскости осуществляет фильтрацию пространственных частот: пропускаются только те гармоники, фурье-образ которых попадает в апертуру щели.

Для наблюдения мультиплицирования поменяем местами сетку $C$ и щель $D$ (рис.~3). Сначала получим резкое изображение щели на экране, затем поместим кассету с сетками в плоскость $F$. При правильной подборке ширины входной щели $D$ и периода сетки на экране наблюдается многократное повторение изображения щели --- мультиплицирование. Уменьшение ширины щели приводит к расширению её фурье-образа, что облегчает ``рассечение'' спектра сеткой и увеличивает число повторений изображения.

\section{Результаты и обсуждения}

\subsection{Сравнение методов определения периода решёток}

В ходе работы периоды шести сеток были определены тремя независимыми методами: по пространственному спектру (дифракция Фраунгофера), по увеличенному изображению в модели микроскопа и по разрешающей способности (метод Аббе). Для наглядного сравнения сведем полученные значения в единую таблицу~\ref{tab:comparison_final}.

\begin{table}[h!]
\centering
\begin{tabular}{|c|c|c|c|}
\hline
№ сетки & $d_{\text{спектр}}$, мкм & $d_{\text{микро}}$, мкм & $d_{\text{Аббе}}$, мкм \\
\hline
1 & $9{,}875 \pm 0{,}101$ & $10{,}95 \pm 0{,}016$ & --- \\
\hline
2 & $24{,}94 \pm 0{,}025$ & $21{,}13 \pm 0{,}033$ & --- \\
\hline
3 & $50{,}34 \pm 0{,}051$ & $48{,}90 \pm 0{,}108$ & $53{,}94 \pm 0{,}25$ \\
\hline
\end{tabular}
\caption{Таблица 6. Сравнение периодов решёток, определённых разными методами}
\label{tab:comparison_final}
\end{table}

Анализ таблицы~\ref{tab:comparison_final} показывает следующее:
\begin{itemize}
    \item Для сетки №\,1 периоды, определённые по спектру ($d_{\text{спектр}} = 9{,}875 \pm 0{,}101$~мкм) и по увеличенному изображению ($d_{\text{микро}} = 10{,}95 \pm 0{,}016$~мкм), расходятся на $\delta \approx 10{,}9\,\%$.

    Это расхождение превышает оценённые случайные погрешности каждого метода, что указывает на наличие систематических ошибок. Наиболее вероятные причины: неточная центровка линз при сборке модели микроскопа, погрешность определения положения промежуточного изображения, а также возможное отклонение от условия малых углов при расчёте по формуле~\eqref{eq:d_spectrum}.
    
    \item Для сетки №\,2 расхождение между спектральным методом и методом микроскопа составляет $\sim 15\,\%$, что также превышает оценённые случайные погрешности. Вероятная причина --- систематическая ошибка при измерении геометрических параметров $a_1$, $b_1$, $a_2$, $b_2$ (неточная центровка линз, погрешность определения положения промежуточного изображения).
    
    \item Для сетки №\,3 все три метода дают значения периода, согласующиеся в пределах $5$--$7\,\%$. Наибольшее расхождение наблюдается между методом Аббе и спектральным методом, что может быть связано с субъективностью определения ``минимального видимого изображения'' при измерении ширины щели~$D$.
\end{itemize}

\subsection{Анализ погрешностей и источников неопределённости}

Основные вклады в суммарную погрешность измерений:
\begin{enumerate}
    \item \textbf{Инструментальная погрешность линейки} ($\sigma_{\text{сист}} = 0{,}5$~мм) --- доминирующий источник для всех методов, основанных на линейных измерениях. Для спектрального метода это приводит к относительной погрешности $\varepsilon_d \approx 1\,\%$ (табл.~\ref{tab:periods_with_errors}).
    
    \item \textbf{Субъективный критерий ``видимости''} в методе Аббе --- определение минимальной ширины щели $D$, при которой ``ещё видно изображение'', зависит от наблюдателя и условий освещения. Это вносит неучтённую систематическую погрешность, которая может достигать $5$--$10\,\%$.
    
    \item \textbf{Неточность установки щели в фокальной плоскости} $F$ --- если щель смещена относительно задней фокальной плоскости линзы Л$_1$, формула~\eqref{eq:d_abbe_final} перестаёт быть строго справедливой, так как изменяется эффективный апертурный угол.
    
    \item \textbf{Приближение малых углов} в формуле~\eqref{eq:d_spectrum} --- для сеток с малым периодом и больших порядков $m$ угол $\theta$ может достигать $10^\circ$--$15^\circ$, при которых $\sin\theta \approx \tan\theta$ даёт погрешность до $2$--$3\,\%$.
    
    \item \textbf{Погрешность фокусных расстояний} линз --- значения $F_1 = 110$~мм и $F_2 = 25$~мм указаны с точностью $\sim 5\,\%$, что напрямую влияет на расчёт увеличения $\Gamma$ и периода по методу Аббе.
\end{enumerate}

\section{Выводы}

\begin{enumerate}
    \item \textbf{По пространственному спектру}: измерены расстояния между дифракционными максимумами для шести сеток; периоды решёток рассчитаны по формуле Фраунгофера с относительной погрешностью $\sim 1\,\%$ (табл.~\ref{tab:periods_with_errors}).
    
    \item \textbf{По увеличенному изображению}: собрана модель проекционного микроскопа, определено увеличение $\Gamma = 6{,}89 \pm 0{,}07$; периоды трёх сеток рассчитаны через геометрическое увеличение, погрешность составила $1{,}5$--$2{,}2\,\%$ (табл.~\ref{tab:periods_microscope}).
    
    \item \textbf{По методу Аббе}: для сетки №\,3 измерена минимальная ширина щели $D = 2{,}17$~мм, период рассчитан по формуле $d = 2F_1\lambda/D$ с погрешностью $0{,}46\,\%$; для остальных сеток измерения ограничены максимальным ходом микрометрического винта.
    
    \item \textbf{Сравнение методов}: значения периода, полученные тремя способами, согласуются в пределах $5$--$15\,\%$; основные расхождения объясняются систематическими погрешностями геометрических измерений и субъективностью визуальных критериев.
    
    \item \textbf{Качественные явления}: продемонстрированы пространственная фильтрация (выделение ориентированных структур через селекцию фурье-компонент) и мультиплицирование (повторение изображения при дискретизации спектра), подтверждающие корпускулярно-волновой дуализм формирования оптических изображений.
\end{enumerate}

Таким образом, в работе экспериментально подтверждена дифракционная природа предела разрешения микроскопа и справедливость формулы Аббе для когерентно освещённых периодических объектов в пределах оценённой экспериментальной погрешности.

\end{document}